\chapter{Abbildungen, die die Struktur erhalten}
\label{b45guide:ch:homomorphisms}

\section{Erhaltung der Operationen}

\begin{definition}[Homomorphismus, Einbettung und Isomorphismus]
\label{b45guide:def:hom}
Eine Abbildung $f:S\to T$ zwischen Halbverbänden heißt
\emph{Halbverbandshomomorphismus}, wenn
\[
 f(x\vee_S y)=f(x)\vee_T f(y)\qquad(x,y\in S).
\]
Eine Abbildung $f:L\to M$ zwischen Verbänden heißt
\emph{Verbandshomomorphismus}, wenn sie beide Operationen erhält:
\[
 f(x\wedge_L y)=f(x)\wedge_M f(y),\qquad
 f(x\vee_L y)=f(x)\vee_M f(y).
\]
Eine \emph{Einbettung} ist ein injektiver Homomorphismus, ein
\emph{Isomorphismus} ein bijektiver Homomorphismus.
Ein \emph{Automorphismus} ist ein Isomorphismus einer Struktur auf sich.
Bei Strukturen mit benannter Null oder Eins wird deren Erhaltung
zusätzlich verlangt, wenn von $0$- oder $(0,1)$-Homomorphismen
gesprochen wird.
\end{definition}

\begin{theorem}[Termtreue und Bestimmtheit durch Erzeuger]
\label{b45guide:thm:term-preservation}
Homomorphismen erhalten die Auswertung aller Terme in den
jeweiligen Operationen. Für einen Verbandsterm $t$ gilt etwa
\[
 f\bigl(t^L(a_1,\ldots,a_n)\bigr)
   =t^M\bigl(f(a_1),\ldots,f(a_n)\bigr).
\]
Zwei Homomorphismen mit derselben Quelle und demselben Ziel,
die auf einer Erzeugermenge übereinstimmen, sind gleich.
Bei benannten Konstanten gilt dies für Terme und Erzeugung
einschließlich dieser Konstanten.
\end{theorem}

\begin{theorem}[Bilder, Urbilder und erzeugte Teilstrukturen]
\label{b45guide:thm:hom-substructures}
Sei $f:S\to T$ ein Halbverbandshomomorphismus oder
$f:L\to M$ ein Verbandshomomorphismus.
Das Bild einer Unterstruktur und das Urbild einer Unterstruktur
sind wieder Unterstrukturen der entsprechenden Art.
Insbesondere ist das gesamte Bild ein Unterhalbverband
beziehungsweise Unterverband des Ziels. Für $X$ in der Quelle gilt
\[
 f[\langle X\rangle]=\langle f[X]\rangle .
\]
Die Erzeugung bezieht sich auf die jeweils erhaltenen Operationen
und gegebenenfalls Konstanten. \cite{supclosed}
\end{theorem}

\begin{theorem}[Isomorphiekriterium]
\label{b45guide:thm:inverse}
Die Umkehrabbildung eines bijektiven Halbverbands- beziehungsweise
Verbandshomomorphismus erhält dieselben Operationen.
Für einen bijektiven Verbandshomomorphismus $f:L\to M$ gilt also
\[
\begin{aligned}
 f^{-1}(u\wedge_M v)&=f^{-1}(u)\wedge_L f^{-1}(v),\\
 f^{-1}(u\vee_M v)&=f^{-1}(u)\vee_L f^{-1}(v).
\end{aligned}
\]
Dasselbe gilt für die Erhaltung benannter Schranken.
Damit sind zwei Strukturen genau dann isomorph, wenn eine
Bijektion zwischen ihnen ihre benannten Operationen erhält.
\end{theorem}

\section{Struktur und induzierte Ordnung}

\begin{theorem}[Ordnungserhaltung und Ordnungsreflexion]
\label{b45guide:thm:order-embedding}
Jeder Halbverbands- oder Verbandshomomorphismus $f$ erfüllt
\[
 x\leq y\quad\Longrightarrow\quad f(x)\leq f(y)
\]
für die zugehörigen induzierten Ordnungen.
Ist $f$ eine Einbettung, so gilt die Äquivalenz
\[
 x\leq y\quad\Longleftrightarrow\quad f(x)\leq f(y).
\]
Eine Bijektion zwischen Halbverbänden gleicher Lesart ist genau
dann ein Halbverbandsisomorphismus, wenn sie den induzierten
Vergleich in beiden Richtungen erhält. Für Verbände gilt
entsprechend: Isomorphie der induzierten Ordnungen und Isomorphie
der beiden Verbandsoperationen sind gleichbedeutend.
\end{theorem}

\begin{theorem}[Wann Schranken automatisch erhalten werden]
\label{b45guide:thm:surjective-bounds}
Ein surjektiver Halbverbandshomomorphismus zwischen Halbverbänden
mit vorhandener Null erhält diese Null; bei vorhandenen Einsen
erhält er auch die Eins. Ein surjektiver Verbandshomomorphismus
zwischen beschränkten Verbänden erhält daher beide Schranken.
Für einen beliebigen Homomorphismus folgt die Erhaltung der
Schranken aus der bloßen Erhaltung der binären Operationen
im Allgemeinen nicht.
\end{theorem}

\begin{definition}[Antiisomorphismus eines Verbandes]
Eine Bijektion $f:L\to M$ heißt \emph{Verbandsantiisomorphismus},
wenn sie die beiden Operationen vertauscht:
\[
 f(x\wedge y)=f(x)\vee f(y),\qquad
 f(x\vee y)=f(x)\wedge f(y).
\]
\end{definition}

\begin{theorem}[Erkennen von Antiisomorphismen]
Eine Bijektion $f:L\to M$ zwischen Verbänden ist genau dann
ein Verbandsantiisomorphismus, wenn
\[
 x\leq y\quad\Longleftrightarrow\quad f(y)\leq f(x).
\]
Die Umkehrabbildung ist ebenfalls ein Antiisomorphismus.
Besitzen beide Verbände Schranken, so werden Null und Eins vertauscht.
\end{theorem}
